\documentclass{article}
\usepackage[utf8]{inputenc}
\usepackage[T1]{fontenc}
\usepackage[english]{babel}
\usepackage{amsmath,amssymb,amsthm}
\usepackage{graphicx}
\usepackage{listings}
\usepackage{tikz}
\usetikzlibrary{matrix,arrows.meta,positioning,calc,fit}
\usepackage{tikz-cd}
\usepackage{hyperref}

\theoremstyle{definition}
\newtheorem{definition}{Definition}[section]
\newtheorem{example}[definition]{Example}
\newtheorem{remark}[definition]{Remark}

\theoremstyle{plain}
\newtheorem{theorem}[definition]{Theorem}

\lstset{
  basicstyle=\tiny\ttfamily,
  breaklines=true,
  columns=fullflexible,
  frame=single,
  showstringspaces=false,
  keepspaces=true
}

\title{AXIO/1:\\
Internal Languages, Categorical Models, and the\\
Spectral Lattice of Formal Systems}
\author{Groupoid Infinity}
\date{\today}

\begin{document}

\maketitle

\begin{abstract}
We analyze the Groupoid Infinity stack of formal languages as a \emph{lattice} of
internal languages, not a linear tower. The languages fall into three distinct
branches: a \emph{constructive type theory and set analysis} spine (Henk--Frank--Christine),
an Homotopy Type Theory and Super Differential Smooth Quantum Spectra \emph{super homotopy} spine (Henk--Per--Anders--Urs),
and a \emph{infinity category theory} spine (Dan--Mike--Ulrik). These spines are connected by functorial relationships
but their syntactic term structures are fundamentally incompatible and should not be
unified into a single AST. In particular, simplicial type theory (Ulrik/STT)
is connected to Urs via directed-homotopy foundations, providing the simplicial
backbone for its 3-layer architecture: Super-Smooth, Differential K-Theory, and Quantum Linear types.
We construct syntactic CW-complexes and chain complexes for each language individually,
define their cohomology groups, and assemble a spectral sequence over the full lattice of language categories.
\end{abstract}

\newpage
\tableofcontents
\newpage

\section{Categorical Foundations}

\begin{definition}[Categories with Families (CwF)]
A CwF is a category $\mathcal{C}$ of contexts and substitutions equipped with
presheaves $Ty, Tm : \mathcal{C}^{\mathrm{op}} \to \mathbf{Set}$ and a
representable natural transformation $p : Tm \to Ty$ satisfying context
comprehension via pullback squares.
\end{definition}

\begin{definition}[Natural Models]
Following Awodey, a natural model is a category $\mathcal{C}$ with a terminal
object $\mathbf{1}$ and a representable natural transformation $p : Tm \to Ty$.
Pullbacks of $p$ along Yoneda provide context extension and substitution.
\end{definition}

\begin{definition}[Grothendieck Fibration]
A functor $p : \mathcal{E} \to \mathcal{B}$ is a fibration when for every
$X \in \mathcal{E}$ and $u : J \to p(X)$ in $\mathcal{B}$ there exists a
Cartesian lift $f : Y \to X$ in $\mathcal{E}$.
\end{definition}

The category of languages $\mathcal{L}$ has formal languages as objects and
syntactic functors (compilers, type checkers, normalizers, extractors) as
morphisms.

\section{The Three Spines of the Language Lattice}
The Groupoid Infinity stack does \emph{not} form a linear tower. It decomposes
into three distinct vertical columns: a constructive/inductive spine on the left
(\textbf{Frank}--\textbf{Christine}--\textbf{Laurent}), a homotopical/cohesive spine in the center
(\textbf{Henk}--\textbf{Per}--\textbf{Anders}--\textbf{Urs}), and an operational/directed branch
on the right (\textbf{Dan}--\textbf{Mike}--\textbf{Ulrik}). The relationships between columns
are functorial (lifts, free constructions, forgetful maps, smooth completions) but the syntaxes
themselves are disjoint.

\begin{center}
\begin{tikzcd}[row sep=large, column sep=large]
& \mathbf{Henk} \ar[d, "\text{+$\Sigma$,Id,W}"] & \\
\mathbf{Frank} \ar[d, "\text{+$\Sigma$,Id}"'] & \mathbf{Per} \ar[d] \ar[l, "\text{+Ind}"] & \mathbf{Dan} \ar[d, dashed, "\text{Lift}"] \\
\mathbf{Christine} \ar[d, "\text{Analysis}"'] \ar[ur, dotted, <->, "\simeq"] \ar[r] & \mathbf{Anders} \ar[d, "\text{Modal}"] \ar[r, dashed, bend left=15, "\text{Forget}"] & \mathbf{Mike} \ar[d, dashed, "\text{Free}"] \\
\mathbf{Laurent} \ar[r, "\text{Smooth}"] & \mathbf{Felix} \ar[d, "\text{Diff}"] & \mathbf{Ulrik} \ar[ul, dashed, "\text{Core}"'] \ar[l, red, bend left=15, "\text{STT}"] \\
& \mathbf{Jack} \ar[d, "\text{Urs}"] & \\
& \mathbf{Urs} & 
\end{tikzcd}
\end{center}

\begin{remark}[Christine $\simeq$ Per]
Christine (full CIC with $\Sigma$, Id, Ind, Prop) and Per (MLTT-80 with
$\Sigma$, Id, W, 0, 1, 2) are at a comparable level of expressiveness. Both
serve as the foundation for Anders and can be considered interchangeable
base systems. The choice between inductive schemes (Frank/Christine path) and
W-types (Per path) is a design decision, not a strict ordering.
\end{remark}

\subsection{The Topoi Structure}
To analyze the gradual expansion of language primitives, we construct a fine-grained lattice of semantic categories radiating from the category of dependent products ($\mathbf{\Pi}\text{-}\mathbf{Cat}$ (Henk)). The system expands in four cardinal directions (adding dependent sums to yield locally cartesian closed categories $\mathbf{LCCC}$, inductive algebras $\mathbf{\Pi}\text{-}\mathbf{Cat} + \mathbf{Ind}$ (Frank), path objects for identity types $\mathbf{\Pi}\text{-}\mathbf{Cat} + \mathbf{Path}$, and bicartesian closed structures $\mathbf{BiCCC}$ (IPL)) before merging into rich categories of models such as $\mathbf{LCCC} + \mathbf{Ind}$ (Christine), $\mathbf{LCCC} + \mathbf{W}$ (Per), $\infty\text{-}\mathbf{Grpd}$ (Anders), and cohesive $\infty$-topoi. On the right, the lattice incorporates the operational algebra branch ($\mathbf{Pres}$ (Dan)), the directed 1-category branch ($\mathbf{1\text{-}Cat}$ (Mike)), and the simplicial $(\infty, 1)$-category branch ($\mathbf{(\infty, 1)\text{-}Cat}$ (Ulrik)). All cross-branch connections are detailed in the global map in Section 2.

\begin{theorem}[Universal Primitive Lattice]
\begin{center}
\scriptsize
\begin{tikzcd}[row sep=large, column sep=0.0em]
& & \mathbf{LCCC} + \mathbf{W}\ \text{(Per)} & \mathbf{LCCC} \ar[l, dashed] & \mathbf{Pres}\ \text{(Dan)} \ar[d, dashed, "\text{Lift}"] \\
\infty\text{-}\mathbf{Grpd}\ \text{(Anders)} & \mathbf{\Pi}\text{-}\mathbf{Cat} + \mathbf{Path} \ar[l, dashed] & \mathbf{\Pi}\text{-}\mathbf{Cat}\ \text{(Henk)} \ar[ur] \ar[dr] \ar[l] \ar[r] & \mathbf{BiCCC}\ \text{(IPL)} \ar[r, dashed] & \mathbf{1\text{-}Cat}\ \text{(Mike)} \ar[d, dashed, "\text{Free}"] \\
\mathbf{Cohesive}\ \infty\text{-}\mathbf{Topos}\ \text{(Felix)} \ar[u, leftarrow, "\text{Modal}"'] & & \mathbf{LCCC} + \mathbf{Ind}\ \text{(Christine)} & \mathbf{\Pi}\text{-}\mathbf{Cat} + \mathbf{Ind}\ \text{(Frank)} \ar[l, dashed] & \mathbf{(\infty, 1)\text{-}Cat}\ \text{(Ulrik)} \\
\mathbf{Diff}\ \infty\text{-}\mathbf{Topos}\ \text{(Jack)} \ar[u, leftarrow, "\text{Diff}"'] & & & & \\
\mathbf{Urs}\ \infty\text{-}\mathbf{Topos}\ \text{(Urs)} \ar[u, leftarrow, "\text{Urs}"'] & & & &
\end{tikzcd}
\end{center}
\end{theorem}

\newpage
\section{Spine I: Set Analysis}
This branch represents the constructive/inductive spine of type theories and functional analysis. It is characterized by inductive algebraic constructions and real/functional analysis foundations.

\subsection{Frank: Minimal CIC}
Extends Henk with strictly positive inductive schemes.
\begin{lstlisting}[language=Caml]
type term =
    | Var of string | Universe of int
    | Pi of string * term * term | Lam of string * term * term | App of term * term
    | Inductive of inductive | Constr of int * inductive * term list
    | Ind of inductive * term * term list * term
and inductive = { name : string; params : (string * term * term) list;
    level : int; constrs : (int * term) list }
\end{lstlisting}

\subsection{Christine: Full CIC (CwF/CwA)}
Complete Calculus of Inductive Constructions with $\Sigma$, Id, and Ind.
\begin{lstlisting}[language=Caml]
type term =
    | Var of name | Universe of level
    | Pi of name * term * term | Lam of name * term * term | App of term * term
    | Sigma of name * term * term | Pair of term * term | Fst of term | Snd of term
    | Id of term * term * term | Refl of term
    | J of term * term * term * term * term * term
    | Inductive of inductive | Constr of int * inductive * term list
    | Ind of inductive * term * term list * term
and inductive = { name : string; params : (name * term * term) list;
    level : level; constrs : (int * term) list }
\end{lstlisting}

\subsection{Laurent: Calculus and Functional Analysis}
Laurent formalizes classical real and functional analysis, distribution theory, and measure theory. It provides first-class support for reals, complex numbers, measures, limits, and integration.
\begin{lstlisting}[language=Caml]
type exp_laurent =
  | Real | Complex | Nat | Bool
  | ConstReal of float | ConstComplex of float * float
  | Forall of string * exp_laurent * exp_laurent
  | Exists of string * exp_laurent * exp_laurent
  | Set of exp_laurent | Member of exp_laurent * exp_laurent
  | Measure of exp_laurent | Integral of exp_laurent * exp_laurent * exp_laurent
  | Seq of (int -> exp_laurent) | Lim of exp_laurent | Inf of exp_laurent | Sup of exp_laurent
\end{lstlisting}

\newpage
\section{Spine II: Super Homotopy Type Theory}
This branch represents the central path of homotopical, cubical, and cohesive type systems, transitioning from pure dependent products to higher-dimensional geometry.

\subsection{Henk: Pure Type Systems (PTS)}
The initial object in LCC categories. Contains only $\Pi$ and universes.
\begin{lstlisting}[language=Caml]
type term =
    | Var of string
    | Universe of int
    | Pi of string * term * term
    | Lam of string * term * term
    | App of term * term
\end{lstlisting}

\subsection{Per: Martin-L\"of Type Theory (MLTT-80)}
MLTT-80 with W-types, base types 0, 1, 2, and cubical primitives.
\begin{lstlisting}[language=Caml]
type exp =
  | EPre of Z.t | EKan of Z.t
  | EVar of ident | EHole
  | EPi of exp * (ident * exp) | ELam of exp * (ident * exp) | EApp of exp * exp
  | ESig of exp * (ident * exp) | EPair of tag * exp * exp
  | EFst of exp | ESnd of exp | EField of exp * string
  | EId of exp | ERef of exp | EJ of exp
  | EPathP of exp | EPLam of exp | EAppFormula of exp * exp
  | EI | EDir of dir | EAnd of exp * exp | EOr of exp * exp | ENeg of exp
  | ETransp of exp * exp | EHComp of exp * exp * exp * exp
  | EPartial of exp | EPartialP of exp * exp | ESystem of exp System.t
  | ESub of exp * exp * exp | EInc of exp * exp | EOuc of exp
  | EEmpty | EIndEmpty of exp
  | EUnit | EStar | EIndUnit of exp
  | EBool | EFalse | ETrue | EIndBool of exp
  | EW of exp * (ident * exp) | ESup of exp * exp | EIndW of exp * exp * exp
\end{lstlisting}

\subsection{Anders: Modal Cubical HoTT}
CCHM cubical structure with Path, Glue, shape modality ($\Im$), and
higher inductive types (Coequ, Disc, Im, Fla).
\begin{lstlisting}[language=Caml]
type value =
  | VKan of Z.t | VPre of Z.t | Var of ident * value | VHole
  | VPi of value * clos | VLam of value * clos | VApp of value * value
  | VSig of value * clos | VPair of tag * value * value | VFst of value | VSnd of value
  | VId of value | VRef of value | VJ of value
  | VPathP of value | VPLam of value | VAppFormula of value * value
  | VI | VDir of dir | VAnd of value * value | VOr of value * value | VNeg of value
  | VTransp of value * value | VHComp of value * value * value * value
  | VGlue of value | VGlueElem of value * value * value | VUnglue of value * value * value
  | VEmpty | VIndEmpty of value | VUnit | VStar | VIndUnit of value
  | VBool | VFalse | VTrue | VIndBool of value
  | W of value * clos | VSup of value * value | VIndW of value * value * value
  | VNat | VZero | VSucc of value | VIndNat of value * value * value
  | VCoequ of value * value * value * value
  | VIm of value | VInf of value | VIndIm of value * value | VJoin of value
  | VFla of value | VFlaUnit of value | VFlaCounit of value | VIndFla of value * value
  ...
and clos = ident * (value -> value)
\end{lstlisting}

\subsection{Urs: The Family of Layered Sublanguages}
Urs is not a single monolithic type system, but rather a family of three
layered sublanguages where each layer builds on top of the syntax, modalities, and constructors of the previous one:

\begin{center}
\resizebox{0.95\textwidth}{!}{%
\begin{tabular}{|c|l|l|}
\hline
\textbf{Sublanguage} & \textbf{Modules} & \textbf{Primitives} \\
\hline
\textbf{I. Felix} & \texttt{graded, flat, sharp, bose,} &
    Bose/Fermi grades, $\flat$, $\sharp$, \\
 & \texttt{fermi, tensor, smth, supsmth} &
    $A \otimes B$, SmthSet, SupSmthSet \\
\hline
\textbf{II. Jack} & \texttt{spectra, grpd, form,} &
    Spectrum, Susp, Grpd, $\Omega^n(X)$, \\
 & \texttt{ku, diff, diffku} &
    $KU^{\tau}_G$, $d$, DiffKU \\
\hline
\textbf{III. Urs} & \texttt{conf, braid, qubit, linear} &
    Config$^n$, Braid$_n$, Qubit, $!A$ \\
\hline
\end{tabular}%
}
\end{center}

\begin{remark}[The 3-Layer Architecture]
The sublanguages represent a progressive structure for formalizing mathematical and
physical concepts:
\begin{enumerate}
  \item \textbf{Felix}: Graded universes and tensors, group
    actions, and super-modality operators for geometric cohesion and
    differential structure---bosonic/fermionic grades, flat/sharp modalities.
  \item \textbf{Jack}: Homotopical foundations via
    spectra and groupoids, supporting differential K-theory with differential
    forms, connections, and refined cohomology structures.
  \item \textbf{Urs}: Configuration spaces and braids for Urs
    statistics, supporting Urs systems and linear types ($!A$) for
    resource-sensitive Urs programming.
\end{enumerate}
Each sublanguage culminates in a system tailored to formalize superpoints
($\mathbb{R}^{m|n}$), supersymmetry, and equivariant structures.
\end{remark}

\subsubsection{Felix (Layer I)}
Extends the homotopy core with graded universes, tensors, group actions, and cohesive modalities ($\flat$, $\sharp$, $\Im$) for synthetic supergeometry.
\begin{lstlisting}[language=Caml]
type grade = Bose | Fermi
type exp_coh =
  | Core of exp_hom
  | UniverseCoh of int * grade
  | SmthSet | Plot of int * exp_coh * exp_coh
  | Flat of exp_coh | Sharp of exp_coh | Shape of exp_coh
  | Bosonic of exp_coh | Tensor of exp_coh * exp_coh
  | SupSmthSet
\end{lstlisting}

\subsubsection{Jack (Layer II)}
Extends cohesion with stable homotopy (spectra, suspensions, wedges) and differential cohomology (forms, connections, and differential K-theory).
\begin{lstlisting}[language=Caml]
type exp_diff =
  | Coh of exp_coh
  | Grpd of int | Comp of int * exp_diff * exp_diff * exp_diff
  | Spectrum | Susp of exp_diff | Wedge of exp_diff * exp_diff
  | HomSpec of exp_diff * exp_diff | KU_G of exp_diff * exp_diff * exp_diff
  | Forms of int * exp_diff | Diff of int * exp_diff
  | DiffKU_G of exp_diff * exp_diff * exp_diff * exp_diff
\end{lstlisting}

\subsubsection{Urs (Layer III)}
The top layer, adding configuration spaces, braids, and linear resource types ($!A$) for Urs programming.
\begin{lstlisting}[language=Caml]
type exp_Urs =
  | DiffK of exp_diff
  | Qubit of exp_Urs * exp_Urs | Linear of exp_Urs
  | AppQubit of exp_Urs * exp_Urs * exp_Urs
  | FuseQubit of exp_Urs * exp_Urs * exp_Urs
  | Config of int * exp_Urs | Braid of int * exp_Urs
\end{lstlisting}

\begin{remark}[STT $\to$ Urs: The Simplicial Connection]
Simplicial type theory (Ulrik/STT) provides a direct connection to Urs.
The directed interval $\mathbb{I}^{\to}$, extension types, and modal
$\Pi$ in STT supply the $\infty$-categorical framework that Urs's Layer I
(Homotopy) builds upon. By isolating the category of categories $\mathbf{Cat}$ as a reflective subuniverse in simplicial type theory, STT provides the required directed homotopy structures. Concretely, STT's directed interval becomes the path/homotopy backbone, while its modal/shape operations inform Urs's cohesive modalities ($\flat$, $\sharp$, $\Im$) in Layer II\@. This is represented by the $\mathrm{STT}$ arrow from Ulrik to Urs in the lattice diagram, which is \emph{not} simply a factoring through Anders.
\end{remark}

\newpage
\section{Spine III: Infinity Category Theory}
This branch encompasses operational algebra and directed category theories, dealing with presentations of algebraic shapes and directed categorical structures.

\subsection{Dan: Operational Algebra}
Dan operates on a completely different level: it manipulates algebraic
presentations (groups, simplicial sets, chain complexes) operationally.
Its term language has no dependent types, no universes, and no $\Pi$/$\Sigma$.
Instead it works with group words, permutations, matrices, face/degeneracy
maps, and boundary operators.

\begin{lstlisting}[language=Caml]
type type_name =
  | Simplex | Group | Simplicial | Chain | Category | Monoid | Ring | Field
  | PermGroup | MatGroup | FpGroup | PcGroup | Module | Set | GSet
type term =
  | Id of string | Comp of term * term | Inv of term | Pow of term * superscript
  | E | Matrix of int list list | Add of term * term | Mul of term * term
  | Perm of int list | Cycle of int list list
  | PermGroup of term list | MatGroup of term list * int * int
  | FpGroup of string list * term list | PcGroup of string list * term list
  | Face of superscript * term | Deg of superscript * term | Boundary of term
  | Act of term * term | Hom of term * term | RelEq of term * term
type hypothesis =
  | Decl of string list * type_name
  | Equality of string * term * term
  | Mapping of string * term * term
  | Presentation of string * term | Relation of term
  | Action of string * term * term
  | Orbit of term * term * term list | Stabilizer of term * term * term
\end{lstlisting}

\subsection{Mike: Directed 1-Category Type Theory}
A \emph{linear} type theory with a bidirectional type checker enforcing
quadraticality (each category variable occurs exactly once covariantly and
once contravariantly).
\begin{lstlisting}[language=Caml]
type cat = CVar of name | COp of cat | CProd of cat * cat
type m_type =
  | MHom of cat * cat_term * cat_term | MTensor of m_type * m_type
  | MCoend of cat * name * m_type | MEnd of cat * name * m_type
  | MFunc of m_type * m_type | MApp of name * cat_term list
type m_term =
  | MTId of cat_term
  | MTJ of m_type * name * name * name * m_term * cat_term * cat_term * m_term
  | MTJCov of m_type * name * m_term * cat_term * m_term
  | MTJContra of m_type * name * m_term * cat_term * m_term
  | MTTensorIntro of m_term * m_term | MTTensorElim of name * name * m_term * m_term
  | MTCoendIntro of name * name * name * m_term
  | MTCoendElim of name * name * m_term * m_term
  | MTEndIntro of name * m_term
  | MTEndElim of name * name * name * m_term * m_term
  | MTFuncIntro of name * m_type * m_term | MTFuncElim of m_term * m_term
  | MTVar of name
\end{lstlisting}

\subsection{Ulrik: Simplicial \texorpdfstring{$\infty$}{Infinity}-Categories}
Extends Mike's foundation with directed interval $\mathbb{I}$, shape
inclusions ($\varphi \subseteq \psi$), extension types, modal $\Pi$, and
twisted arrow categories ($A^{\mathrm{tw}}$).

To construct a type theory for synthetic category theory, one might hope to interpret type theory in the category of categories ($\infty$-categories or otherwise) to ensure that types realize categories. However, the category $\mathbf{Cat}$ of small categories is too poorly behaved to form a model of Martin-Löf Type Theory (MLTT). Instead, Riehl and Shulman (2017) embed $\mathbf{Cat}$ as a reflective subcategory in the category of ($\infty$-)presheaves on the simplex category $\widehat{\Delta}$, which is sufficiently rich to model HoTT. Simplicial Type Theory (STT) then axiomatizes a portion of $\widehat{\Delta}$ to isolate $\mathbf{Cat}$ as a reflective subuniverse within the type theory.

\begin{lstlisting}[language=Caml]
type exp =
  | EUniv | EVar of name | EPi of exp * (name * exp) | ELam of (name * exp option) * exp
  | EApp of exp * exp | ESig of exp * (name * exp) | EPair of exp * exp
  | EFst of exp | ESnd of exp | EId of exp * exp * exp | ERef of exp
  | EIDir | EZeroDir | EOneDir | ELeq of exp * exp | EShapeInc of exp * exp
  | ESystem of (exp * exp) list | EExt of exp * exp * exp
  | EModalPi of exp * (name * exp) | EModalLam of (name * exp) * exp
  | EModalApp of exp * exp | ETw of exp | ETwPi0 of exp | ETwPi1 of exp
  | EJoin of exp * exp | EMeet of exp * exp | ENeg of exp
  | EJ of exp * name * name * name * exp * exp * exp * exp
  | EJCov of exp * name * exp * exp * exp | EJContra of exp * name * exp * exp * exp
\end{lstlisting}

\newpage
\section{Cohomology of Programming Languages}
Each language's syntax is analyzed \emph{individually} as a CW-complex.

For a formal language $\mathcal{O}_x$, we associate a chain complex of abelian groups (free $\mathbb{Z}$-modules) $C_*(\mathcal{O}_x)$ constructed from the syntax and reduction/conversion semantics. This construction models the syntactic space as a cell complex (CW-complex) where paths and higher homotopies encode the computational rules.

\begin{definition}[Syntactic CW-Complex]
For a language $\mathcal{O}_x$ with term syntax $T$ generated by a signature $\Sigma$, the syntactic cell complex $C_*(\mathcal{O}_x)$ is defined in low dimensions by:
\begin{itemize}
    \item \textbf{0-cells ($C_0(\mathcal{O}_x)$)}: The free $\mathbb{Z}$-module generated by the set of basic syntactic constructors $c \in \Sigma$ (e.g., \texttt{Var}, \texttt{Pi}, \texttt{Lam}, \texttt{App} for Henk). Any term $t \in T$ has an associated constructor-occurrence representation:
    \[
    [t] = \sum_{c \in \Sigma} \mathrm{count}(c, t) \cdot c \;\; \in C_0(\mathcal{O}_x)
    \]
    where $\mathrm{count}(c, t)$ is the number of times the constructor $c$ appears in the syntax tree of $t$.
    
    \item \textbf{1-cells ($C_1(\mathcal{O}_x)$)}: The free $\mathbb{Z}$-module generated by the rewriting and reduction rules $r \in \mathcal{R}$ of the language (e.g., $\beta$-reductions, evaluation steps). Each $1$-cell $r : t \to t'$ represents a directed rewrite path from a redex term $t$ to a contractum term $t'$.
    
    \item \textbf{2-cells ($C_2(\mathcal{O}_x)$)}: The free $\mathbb{Z}$-module generated by equivalence relations, critical pairs, confluence diagrams, and conversion equations (e.g., $\eta$-conversions, commuting conversions). A $2$-cell $e$ represents a coherence condition or a closed loop of rewrite steps.
    
    \item \textbf{Higher cells ($C_k(\mathcal{O}_x)$ for $k \geq 3$)}: Coherence relations between equations, such as Mac Lane pentagon relations, path-type univalence coherences (in Anders), or quadraticality constraints (in Mike).
\end{itemize}
\end{definition}

\begin{definition}[Boundary Operators]
The boundary operators $\partial_k : C_k(\mathcal{O}_x) \to C_{k-1}(\mathcal{O}_x)$ are defined as follows:
\begin{enumerate}
    \item The 1-boundary operator $\partial_1 : C_1(\mathcal{O}_x) \to C_0(\mathcal{O}_x)$ is the linear map determined by the net change in constructor counts caused by a rewrite rule $r : t \to t'$:
    \[
    \partial_1(r) = [t'] - [t]
    \]
    
    \item The 2-boundary operator $\partial_2 : C_2(\mathcal{O}_x) \to C_1(\mathcal{O}_x)$ is determined by the boundary of confluence diagrams or conversion loops. If a 2-cell $e$ represents a confluence diagram from a term $s$ reducing to $s'$ via two different paths of rewrites $p_1 = (r_{1,1}, \dots, r_{1,m})$ and $p_2 = (r_{2,1}, \dots, r_{2,n})$, the boundary is:
    \[
    \partial_2(e) = \sum_{j=1}^m r_{1,j} - \sum_{k=1}^n r_{2,k}
    \]
    Since both paths $p_1$ and $p_2$ start at $s$ and end at $s'$, the net constructor changes must match:
    \[
    \partial_1(p_1) = [s'] - [s] = \partial_1(p_2)
    \]
    Thus:
    \[
    \partial_1(\partial_2(e)) = \partial_1(p_1) - \partial_1(p_2) = 0
    \]
    This proves that $\partial_1 \circ \partial_2 = 0$, establishing a valid chain complex.
\end{enumerate}
\end{definition}

\begin{definition}[Syntactic Cohomology]
Let $C_*(\mathcal{O}_x)$ be the syntactic chain complex of $\mathcal{O}_x$ and let $G$ be an abelian group of coefficients. The cochain complex $C^*(\mathcal{O}_x; G) = \mathrm{Hom}(C_*(\mathcal{O}_x), G)$ has coboundary operators $\delta^k : C^k(\mathcal{O}_x; G) \to C^{k+1}(\mathcal{O}_x; G)$ defined by $\delta^k(\phi) = \phi \circ \partial_{k+1}$.
The syntactic cohomology groups are:
\[
H^k(\mathcal{O}_x; G) = \ker(\delta^k) / \mathrm{im}(\delta^{k-1})
\]
\end{definition}

\subsection{Physical Interpretations of Cohomology Groups}
The cohomology groups $H^k(\mathcal{O}_x; G)$ measure fundamental properties of the syntax and semantics of $\mathcal{O}_x$:
\begin{itemize}
    \item $H^0(\mathcal{O}_x; G)$: Measures the \emph{syntactic dimension} or independent constructor classes. A class in $H^0$ is a function $f : C_0 \to G$ such that for any rewrite $r : t \to t'$, $f([t]) = f([t'])$. Thus, $H^0$ captures invariants preserved under computation (such as type preservation or conserved weights of terms).
    \item $H^1(\mathcal{O}_x; G)$: Measures \emph{normalization obstructions}. An element in $H^1$ is represented by a 1-cocycle (an assignment of weights to rewrites that sums to zero on closed loops) modulo those coming from constructor valuations. Non-trivial $H^1$ classes indicate non-confluent rewrites, cycles of reductions (non-termination), or irreversible compilation steps.
    \item $H^2(\mathcal{O}_x; G)$: Measures \emph{coherence obstructions}. It captures the obstructions to filling confluence diagrams with higher-dimensional conversion cells. A non-trivial class in $H^2$ represents a critical pair (or confluence loop) that cannot be closed by conversions or equivalence cells, pointing to deep semantic mismatches or failures of univalence/coherence.
\end{itemize}

\subsection{Examples across the Branches}

\begin{example}[Branch I/II: Constructive \& Homotopical Spine]
\end{example}
\begin{itemize}
    \item \textbf{Henk ($\mathcal{O}_\Pi$)}: The core type theory.
    \[
    C_0 = \{ \mathtt{Var}, \mathtt{Universe}, \mathtt{Pi}, \mathtt{Lam}, \mathtt{App} \}
    \]
    The 1-cells are generated by $\beta$-reduction:
    \[
    r_\beta : \mathtt{App}(\mathtt{Lam}(x, A, t), u) \to t[u/x]
    \]
    $\partial_1(r_\beta) = [t[u/x]] - ( \mathtt{App} + \mathtt{Lam} + [t] + [u] )$.
    The 2-cells consist of $\eta$-conversions:
    \[
    e_\eta : \mathtt{Lam}(x, A, \mathtt{App}(t, x)) \equiv_\eta t
    \]
    Since Henk is strongly normalizing and confluent (Church-Rosser), the higher cohomology $H^k(\mathcal{O}_\Pi; \mathbb{Z})$ for $k \geq 1$ is trivial, reflecting a contractible syntactic space.
    \item \textbf{Anders ($\mathcal{O}_{\mathrm{Anders}}$)}: Adds univalence and Glue types. The univalence axiom represents a path-space equivalence between isomorphic types. This introduces higher-dimensional coherences (Glue cells) acting as 3-cells, preventing non-trivial 2-dimensional obstructions.
    \item \textbf{Urs ($\mathcal{O}_{\mathrm{Urs}}$)}: Adds modalities ($\flat, \sharp$) and differential structures (forms, connections). The syntactic CW-complex of $\mathcal{O}_{\mathrm{Urs}}$ is enriched with topological and differential cells, yielding non-trivial $H^k$ corresponding to de Rham cohomology and Chern-Weil invariants of smooth/differential type structures.
\end{itemize}

\begin{example}[Branch III: Operational \& Directed Branch]
\end{example}
\begin{itemize}
    \item \textbf{Dan ($\mathcal{O}_{\mathrm{Dan}}$)}: Operational algebra. Here, term syntax encodes algebraic operations (e.g., face maps $\partial_i$, degeneracies $s_i$). The cell complex $C_*(\mathcal{O}_{\mathrm{Dan}})$ is isomorphic to the standard simplicial chain complex of the underlying category/algebraic presentation.
    \item \textbf{Mike ($\mathcal{O}_{\mathrm{Mike}}$)}: Directed category theory. 1-cells are non-reversible functors/morphisms, meaning paths are directed. The boundary operator $\partial_1$ tracks directed change. Quadraticality constraints act as 2-cells that model the non-commutative relation profiles of monoidal/directed types.
    \item \textbf{Ulrik ($\mathcal{O}_{\mathrm{Ulrik}}$)}: Simplicial $\infty$-categories. Generalizes Mike's 1-categories to directed simplicial spaces. The higher cells ($C_k$ for $k \geq 2$) represent the Segal completeness conditions and Rezk completeness conditions, ensuring that the directed paths compose coherently up to higher homotopies.
\end{itemize}

\newpage
\section{Spectral Sequence of Language Categories}
Unlike a simple linear filtration, the language category $\mathcal{L}$ is
organized as a \emph{lattice} with multiple join-paths. The spectral sequence
is defined over the partial order of syntactic subcategories:
\[
E_1^{p,q} = H^q(\mathcal{O}_p) \implies H^{p+q}(\mathcal{L})
\]

\begin{table}[ht]
\centering
\resizebox{\textwidth}{!}{%
\begin{tabular}{c|l|l|l|l}
\hline
\textbf{$p$} & \textbf{Language} & \textbf{Branch} & \textbf{Primitives added}
& \textbf{$|C_0|$} \\
\hline
0 & Henk & CTT & $\Pi$, $U^n$ & 5 \\
1 & Frank & CTT & + Ind & 8 \\
2a & Christine & CTT & + $\Sigma$, Id & 14 \\
2b & Per & CTT & + $\Sigma$, Id, W, 0, 1, 2 & 26 \\
3 & Anders & CTT & + Path, I, Glue, HComp, HIT & 40+ \\
4a & Felix & CTT/Modal & + $\flat$, $\sharp$, $\Im$, $\otimes$, Plot & 10 \\
4b & Jack & CTT/Modal & + Spectra, Susp, Grpd, Forms, DiffKU & 12 \\
4c & Urs & CTT/Modal & + Qubit, Linear, Config, Braid & 7 \\
\hline
--- & Dan & Algebra & Groups, Simplices, Chains, Boundaries & 20 \\
--- & Mike & Directed & Hom, $\otimes$, $\int$, $\prod$, $\multimap$ & 13 \\
--- & Ulrik & Directed & + $\mathbb{I}^{\to}$, Ext, Tw, $\leq$ & 25 \\
\hline
\end{tabular}%
}
\caption{The language lattice: filtration indices and constructor counts.
Note the split at $p=2$ (Christine/Per), the four sublanguages of Urs at $p=4$ (Homotopy, SupSmth, DiffK, Urs), and the separate branch indices
for Dan, Mike, and Ulrik.}
\end{table}

\begin{definition}[Morphisms in $\mathcal{L}$]
Morphisms $f : \mathcal{O}_x \to \mathcal{O}_y$ are classified as:
\begin{enumerate}
    \item $\mathcal{F}_{\mathrm{enrich}}$: Enrichment (type inference, elaboration).
    \item $\mathcal{F}_{\mathrm{simp}}$: Simplification (type erasure, extraction).
    \item $\mathcal{F}_{\mathrm{trans}}$: Translation (compilation across branches).
    \item $\mathcal{F}_{\mathrm{eval}}$: Evaluation (normalization to values).
    \item $\mathcal{F}_{\mathrm{norm}}$: Normalization ($\beta\eta$-reduction).
    \item $\mathcal{F}_{\mathrm{certify}}$: Certification (proof checking).
\end{enumerate}
Cross-branch morphisms (e.g.\ \textbf{Lift}: Dan $\to$ Mike, \textbf{Core}:
Ulrik $\to$ Anders) are structure-changing and do not preserve the term type.
\end{definition}

\subsection{Functorial Conversions and Cross-Branch Morphisms}
The language category $\mathcal{L}$ is not merely a set of objects; the relationships between the branches are mediated by structured functorial conversions. The primary conversions, their statuses, and associated information losses are detailed in Table~\ref{tab:functors}.

\begin{table}[ht]
\centering
\resizebox{\textwidth}{!}{%
\begin{tabular}{|l|c|c|l|p{8cm}|}
\hline
\textbf{Functor / Conversion} & \textbf{Source} & \textbf{Target} & \textbf{Status} & \textbf{Semantics \& Information Loss} \\
\hline
\textbf{Lift} & Dan & Mike & Yes & Semantical lifting of presentation generators and equations to 1-categorical structures. \\
\hline
\textbf{Forget} & Mike & Dan & Partial & Truncates categories back to operation-based presentation generators. \\
\hline
\textbf{Free} & Mike & Ulrik & Yes & Fully faithful embedding of directed 1-categories into directed $\infty$-categories. \\
\hline
\textbf{Trunc} & Ulrik & Mike & No & Cannot be done functorially without losing higher-dimensional coherence. \\
\hline
\textbf{Core} & Ulrik & Anders & Yes & Groupoid Core: discards directed path structure, retaining only the sub-groupoid of isomorphisms. \\
\hline
\textbf{Disc} & Anders & Ulrik & Yes & Discrete embedding: interprets $\infty$-groupoids as discrete $\infty$-categories via disc modality. \\
\hline
\textbf{Forget} & Anders & Mike & Partial & 1-category approximation: projects higher paths onto 1-category hom-sets. \\
\hline
\textbf{Groupoid} & Mike & Anders & Yes & Forgetful functor from directed 1-categories to spaces of isomorphism paths. \\
\hline
\textbf{Lift} (Composite) & Dan & Ulrik & Yes & Composition of \textbf{Lift} and \textbf{Free}: maps operations to Rezk $\infty$-categories. \\
\hline
\textbf{Forget} & Ulrik & Dan & No & Directed simplicial spaces are too rich to map back to discrete operations directly. \\
\hline
\textbf{Forget} & Anders & Dan & No & Undirected spaces cannot be functorially mapped to discrete presentation equations. \\
\hline
\textbf{Modal} & Anders & Felix & Yes & Extends cubical HoTT with cohesive modalities ($\flat$, $\sharp$, $\Im$) and super-grades. \\
\hline
\textbf{Diff} & Felix & Jack & Yes & Stable extension to spectra, suspensions, forms, and differential K-theory. \\
\hline
\textbf{Urs} & Jack & Urs & Yes & Extends differential K-theory with configuration spaces, braids, and linear types. \\
\hline
\end{tabular}%
}
\caption{Primary conversions and functors between the language layers.}
\label{tab:functors}
\end{table}

\subsubsection{Deep Dive: The Lift Functor (\texorpdfstring{$\mathbf{Dan} \to \mathbf{Mike} / \mathbf{Ulrik}$}{Dan -> Mike / Ulrik})}
The \textbf{Lift} functor acts as a semantical compilation bridge that transports combinatorial presentations from the operational layer $\mathbf{Dan}$ to the directed synthetic type theories of $\mathbf{Mike}$ and $\mathbf{Ulrik}$. Rather than a simple syntax translator, it elaborates discrete presentations into rich categorical models in four distinct phases:
\begin{enumerate}
    \item \textbf{Signature Ingestion}: Reads the generator set, face maps ($\partial_i$), degeneracy maps, and algebraic presentation relations directly from the $\mathbf{Dan}$ syntax.
    \item \textbf{Semantical Lifting}: Maps the ingested signature into category-theoretic structures, translating generator equations into functors, categories, and sheaves.
    \item \textbf{Elaboration to STT}: Translates composition equations to identity paths ($\mathtt{EId}$) in $\mathbf{Ulrik}$. The directed interval $\mathbb{I}^{\to}$ is employed to model the domains and boundary conditions.
    \item \textbf{Metatheory Absorption}: Automatically equips the lifted object with the abstract category-theoretic library of $\mathbf{Ulrik}$ (including Yoneda lemmas, limits, adjunctions, and Kan extensions).
\end{enumerate}

\subsubsection{Deep Dive: The Groupoid Core (\texorpdfstring{$\mathbf{Ulrik} \to \mathbf{Anders}$}{Ulrik -> Anders})}
The translation from $\mathbf{Ulrik}$ to $\mathbf{Anders}$ is a structural forgetful conversion that extracts the underlying homotopy type (the $\infty$-groupoid of isomorphisms) from a directed $\infty$-category. Since the simplicial directed interval $\mathbb{I}^{\to}$ and directed hom-types ($x \to y$) of $\mathbf{Ulrik}$ do not exist in the cubical undirected world of $\mathbf{Anders}$, the functor operates as follows:
\begin{itemize}
    \item Discards directed arrows, keeping only the sub-groupoid $A^{\mathrm{core}}$ containing morphisms with valid inverses.
    \item Discards the Segal and Rezk completeness conditions, transforming directed category composition into symmetric path composition.
    \item Replaces the directed interval $\mathbb{I}^{\to}$ with the cubical symmetric interval $I$ (with connection and reversal operations).
\end{itemize}

\subsubsection{Case Study: M\"obius Strip Representation}
To illustrate the conversions, Table~\ref{tab:mobius} compares the representation of a M\"obius Strip across the three primary layers ($\mathbf{Dan}$, $\mathbf{Ulrik}$, and $\mathbf{Anders}$).

\begin{table}[ht]
\centering
\resizebox{\textwidth}{!}{%
\begin{tabular}{|p{3cm}|p{4.5cm}|p{4.5cm}|p{4.5cm}|}
\hline
\textbf{Aspect} & \textbf{Dan (Operational)} & \textbf{Ulrik (Simplicial STT)} & \textbf{Anders (Cubical HoTT)} \\
\hline
\textbf{Framework} & Computer algebra presentations & Simplicial Type Theory with $\mathbb{I}^{\to}$ & Cubical HoTT with modal operators \\
\hline
\textbf{Mathematical Type} & Combinatorial simplex complex with boundary relations & Rezk $\infty$-category built from simplicial diagrams & Univalent fiber bundle: $\sum_{x: S^1} \mathrm{Moebius}(x)$ \\
\hline
\textbf{Twist Encoding} & Equations on face maps ($\partial_i$) & Directed paths over $\mathbb{I}^{\to}$ satisfying Segal conditions & Univalent path transport ($\mathtt{ua}$) using boolean negation ($\mathtt{true} \equiv \mathtt{false}$) \\
\hline
\textbf{Homotopy Type} & Contractible local complex ($\ast$) & $\infty$-category whose core homotopy type is $S^1$ & Homotopy equivalence to the circle $S^1$ \\
\hline
\textbf{Abstraction Level} & Combinatorial building blocks & Global synthetic $\infty$-category & Global synthetic univalent bundle \\
\hline
\textbf{Functor Action} & Ingested by \textbf{Lift} & Lifted to directed $\infty$-category & Extracted via \textbf{Core} as undirected $\infty$-groupoid \\
\hline
\textbf{Information Kept} & Full combinatorial generators & Full directed structure & Homotopy skeleton; directed composition is lost \\
\hline
\end{tabular}%
}
\caption{Comparison of the M\"obius Strip representation across the layers.}
\label{tab:mobius}
\end{table}

\end{document}